%% ma: L12-B08-0163
%% lop: 12
%% chuong: 2
%% bai: 8
%% dang: 12.8.52
%% loai: TN4
%% muc: VD
%% dap_an: "B"
%% nguon: "CĐ2.4.pdf giáo viên gửi 27/10/2026 (M07, Lê Minh Khôi – Nguyễn Anh Khôi, Chương II Chủ đề 4 Các bài toán về tâm tỉ cự), Câu 13"
%% kien_thuc: "Bài 8 (biểu thức tọa độ của tích vô hướng), lớp 10 (giá trị nhỏ nhất của tam thức bậc hai)"
%% trang_thai: chinh-thuc
%% ghi_chu: "Giáo viên đã duyệt 27/10/2026"
\begin{ex}
Trong không gian $Oxyz$, cho ba điểm $A(4;1;5)$, $B(3;0;1)$, $C(1;-2;0)$ và $K$ là điểm di động thuộc mặt phẳng $(Oxz)$. Tính giá trị nhỏ nhất của biểu thức $S=\overrightarrow{KA}\cdot\overrightarrow{KB}+\overrightarrow{KB}\cdot\overrightarrow{KC}+\overrightarrow{KC}\cdot\overrightarrow{KA}$.
\choice{$S_{\min}=\dfrac{34}3$}{\True $S_{\min}=-\dfrac{34}3$}{$S_{\min}=\dfrac{35}3$}{$S_{\min}=-\dfrac{35}3$}
\loigiai{$K(x;0;z)$: $S=3x^2-16x+3z^2-12z+22=3\left(x-\dfrac83\right)^2+3(z-2)^2-\dfrac{34}3$.

$S_{\min}=-\dfrac{34}3$ khi $K\left(\dfrac83;0;2\right)$.}
\end{ex}
